\scanpage{53}
\mspage{9}
% archive sheet number: 52
% slanted
\margin{\scriptsize\uline{NB} Cette opération a l'air d'un \uline{affaiblissement} de structure --- mais en affaiblissant, sur une composante $C$ donnée de $\unsure{\Sigma^{+}|\widetilde{\Sigma}}$, lorsqu'on ``gomme'' \uline{tous} les repères de $C$ --- \unsure{donc} c'est au contraire plutôt dans la nature d'un renforcement de structure, --- ou plutôt, du moment qu'on se sert des $R'_C=R'\cap\unsure{C|X}$, dans les catégories topologiques, pour distinguer l'opération de gommage \unsure{et celle de} \deleted{\unsure{enlever tout}} de \unsure{sous-marquage} (cf plus bas), \unsure{à aucun|aucun de} \illegible{2} sont inverses. \unsure{Deux autres|Décrire} des \illegible{2} sont des équivalences de groupoïdes. \unsure{Alors, il faudrait|Il faudrait alors} séparer les deux cas, \unsure{pour distinguer nos} \uline{opérations de gommage essentielles} \unsure{(mieux|moyen)}, et (\unsure{qui change} \illegible{2}) \unsure{pour les} \unsure{cf} les contours.}

\uline{NB} Ayant (via un groupoïde équivalent \deleted{\unsure{de}} à description purement combinatoire) ``reconstruit'' \unsure{ou} \deleted{\unsure{proprement}} réalisé le groupoïde $S{\unsure{\mathscr T|\mathcal T}}D_{\mathrm{top}}^{\pm}$, et le foncteur
\diagram{
\begin{tikzcd}[ampersand replacement=\&]
S\mathscr T D_{\mathrm{top}}^{\pm}\arrow[r]\&(\mathrm{Gr.\ pond.\ rep})
\end{tikzcd}
}
\unsure{on aura} \deleted{\unsure{par|pas}} donc aussi \deleted{\unsure{une formulation de}} \inserted{l'équivalent de} \deleted{\unsure{de}} la donnée d'un $(X,\Sigma,R)$, muni \uline{en plus} d'une partie $R'\subset R$ --- \deleted{\unsure{d'où}} i.e. \unsure{avec} la construction d'un groupoïde \unsure{représentant}
\diagram{
\begin{tikzcd}[ampersand replacement=\&]
S_{*}\mathscr T D_{\mathrm{top}}^{\pm}\arrow[d]\&
\text{(à sous ens. marqué }R'\text{)}
\quad\begin{aligned}&\text{de sommets-}\\&\text{repères}\end{aligned}\\
S_{*}\mathscr T D_{\mathrm{top}}^{\pm}\&
\end{tikzcd}
}
Et \deleted{\unsure{donc}} il nous faudra donc \inserted{décrire} un foncteur ``gommage en $R'$''
\diagram{
\begin{tikzcd}[ampersand replacement=\&]
S_{*}\mathscr T D_{\mathrm{top}}^{\pm}
  (\underset{\text{marqué}}{\overset{\text{sous}}{\deleted{\unsure{R}}R'}})
\arrow[r,"\text{gommage de }R'"]\&{}\arrow[r]\&
S_{*}\mathscr T D_{\mathrm{top}}^{\pm}
\end{tikzcd}
}
\deleted{Cela s'exprime, sur le}
donnant lieu à un diagramme comm. de foncteurs
\diagram{
\begin{tikzcd}[ampersand replacement=\&]
S_{*}\mathscr T D_{\mathrm{top}}^{\pm}(\ldots R'\ldots)
  \arrow[rr,"\text{gommage de }R'"]\arrow[d,"\substack{\text{oubli du marquage}\\R'}"]\&
\&S_{*}\mathscr T D_{\mathrm{top}}^{\pm}\arrow[ddl]\\
S_{*}\mathscr T D\deleted{\unsure{R}}_{\mathrm{top}}^{\pm}\arrow[dr]\&
\&\\
\&(\mathrm{Gra}\deleted{\unsure{\mathrm{b}}}\mathrm{bipond})
\arrow[r]\&\substack{\text{graphes avec}\\\text{pondération}\\\text{par }\unsure{\gamma|\sigma}\text{, mais}\\\text{non par }r\text{,}}
\end{tikzcd}
}
alors qu'au niveau
